\documentclass[10pt,a4paper]{amsart}
\usepackage{amsmath,amssymb,amsthm,graphicx}
\usepackage{verbatim}
\newenvironment{reference}{\comment}{\endcomment}
%\newenvironment{reference}{}{}
\newenvironment{slogan}{\comment}{\endcomment}
\newenvironment{history}{\comment}{\endcomment}
% For commutative diagrams we use Xy-pic
\usepackage[all]{xy}

% We use 2cell for 2-commutative diagrams.
\xyoption{2cell}
% Theorem environments.
%
\theoremstyle{plain}
\newtheorem{theorem}[subsection]{Theorem}
\newtheorem{proposition}[subsection]{Proposition}
\newtheorem{lemma}[subsection]{Lemma}

\theoremstyle{definition}
\newtheorem{definition}[subsection]{Definition}
\newtheorem{example}[subsection]{Example}
\newtheorem{exercise}[subsection]{Exercise}
\newtheorem{situation}[subsection]{Situation}

\theoremstyle{remark}
\newtheorem{remark}[subsection]{Remark}
\newtheorem{remarks}[subsection]{Remarks}

\numberwithin{equation}{subsection}

% Macros
%
\def\lim{\mathop{\mathrm{lim}}\nolimits}
\def\colim{\mathop{\mathrm{colim}}\nolimits}
\def\Spec{\mathop{\mathrm{Spec}}}
\def\Hom{\mathop{\mathrm{Hom}}\nolimits}
\def\Ext{\mathop{\mathrm{Ext}}\nolimits}
\def\SheafHom{\mathop{\mathcal{H}\!\mathit{om}}\nolimits}
\def\SheafExt{\mathop{\mathcal{E}\!\mathit{xt}}\nolimits}
\def\Sch{\mathit{Sch}}
\def\Mor{\mathop{\mathrm{Mor}}\nolimits}
\def\Ob{\mathop{\mathrm{Ob}}\nolimits}
\def\Sh{\mathop{\mathit{Sh}}\nolimits}
\def\NL{\mathop{N\!L}\nolimits}
\def\CH{\mathop{\mathrm{CH}}\nolimits}
\def\proetale{{pro\text{-}\acute{e}tale}}
\def\etale{{\acute{e}tale}}
\def\QCoh{\mathit{QCoh}}
\def\Ker{\mathop{\mathrm{Ker}}}
\def\Im{\mathop{\mathrm{Im}}}
\def\Coker{\mathop{\mathrm{Coker}}}
\def\Coim{\mathop{\mathrm{Coim}}}

% Boxtimes
%
\DeclareMathSymbol{\boxtimes}{\mathbin}{AMSa}{"02}

%
% Macros for moduli stacks/spaces
%
\def\QCohstack{\mathcal{QC}\!\mathit{oh}}
\def\Cohstack{\mathcal{C}\!\mathit{oh}}
\def\Spacesstack{\mathcal{S}\!\mathit{paces}}
\def\Quotfunctor{\mathrm{Quot}}
\def\Hilbfunctor{\mathrm{Hilb}}
\def\Curvesstack{\mathcal{C}\!\mathit{urves}}
\def\Polarizedstack{\mathcal{P}\!\mathit{olarized}}
\def\Complexesstack{\mathcal{C}\!\mathit{omplexes}}
% \Pic is the operator that assigns to X its picard group, usage \Pic(X)
% \Picardstack_{X/B} denotes the Picard stack of X over B
% \Picardfunctor_{X/B} denotes the Picard functor of X over B
\def\Pic{\mathop{\mathrm{Pic}}\nolimits}
\def\Picardstack{\mathcal{P}\!\mathit{ic}}
\def\Picardfunctor{\mathrm{Pic}}
\def\Deformationcategory{\mathcal{D}\!\mathit{ef}}
\usepackage{cleveref}
\crefname{originallemma}{original lemma}{original lemmas}
\Crefname{originallemma}{Original lemma}{Original lemmas}
\makeatletter
\newlabel{lemma-base-change-flat}{{lemma-base-change-flat}{source}{}{sourcepointer-1}{}}
\newlabel{lemma-base-change-flat@cref}{{[originallemma][0][]lemma-base-change-flat}{[1][0][]source}{}{}{}}
\makeatother
\crefname{originalreference}{original reference}{original references}
\Crefname{originalreference}{Original reference}{Original references}
\makeatletter
\newlabel{spaces-morphisms-lemma-base-change-module-flat}{{spaces-morphisms-lemma-base-change-module-flat}{source}{}{sourcepointer-2}{}}
\newlabel{spaces-morphisms-lemma-base-change-module-flat@cref}{{[originalreference][0][]spaces-morphisms-lemma-base-change-module-flat}{[1][0][]source}{}{}{}}
\makeatother
\crefname{originalreference}{original reference}{original references}
\Crefname{originalreference}{Original reference}{Original references}
\makeatletter
\newlabel{spaces-morphisms-lemma-composition-module-flat}{{spaces-morphisms-lemma-composition-module-flat}{source}{}{sourcepointer-3}{}}
\newlabel{spaces-morphisms-lemma-composition-module-flat@cref}{{[originalreference][0][]spaces-morphisms-lemma-composition-module-flat}{[1][0][]source}{}{}{}}
\makeatother
\crefname{originalreference}{original reference}{original references}
\Crefname{originalreference}{Original reference}{Original references}
\makeatletter
\newlabel{spaces-morphisms-lemma-flat-permanence}{{spaces-morphisms-lemma-flat-permanence}{source}{}{sourcepointer-4}{}}
\newlabel{spaces-morphisms-lemma-flat-permanence@cref}{{[originalreference][0][]spaces-morphisms-lemma-flat-permanence}{[1][0][]source}{}{}{}}
\makeatother
\crefname{originalreference}{original reference}{original references}
\Crefname{originalreference}{Original reference}{Original references}
\makeatletter
\newlabel{stacks-properties-lemma-points-cartesian}{{stacks-properties-lemma-points-cartesian}{source}{}{sourcepointer-5}{}}
\newlabel{stacks-properties-lemma-points-cartesian@cref}{{[originalreference][0][]stacks-properties-lemma-points-cartesian}{[1][0][]source}{}{}{}}
\makeatother
\title{Flatness at a point of a morphism of algebraic stacks is characterised by local rings}
\author{The Stacks Project Authors}
\date{stacks-project:stacks-morphisms.tex:local-snapshot-20261012; GNU FDL 1.2 or later}
\begin{document}
\maketitle
\noindent\textit{Editorial note.} All delivered mathematical text is the author's own, copied from the pinned Stacks Project chapter source sources/stacks-algebraic-stacks-morphisms/pinned-original-source.tex; no statement or proof step was rewritten, reordered or removed.\par\medskip
\noindent\textit{Editorial note.} Preamble provenance: the Stacks Project keeps one single shared preamble for every chapter; this card takes it directly from the staged sources/stacks-algebraic-stacks-morphisms/preamble.tex as three verbatim spans (lines 8-12, 14-18 and 154-217), so no derived file was written and nothing was added or redefined. The wrapper class and the base packages are supplied by the assembler.\par\medskip
\noindent\textit{Editorial note.} Bibliography provenance: neither the counted statement nor its proof contains a citation (this lane rescans the whole span before writing the card), so the card carries no bibliography and no entry was reconstructed.\par\medskip
\noindent\textit{Editorial note.} Selection provenance: this card was opened under the file-based control contract (CONTROL.md), with the mechanical pre-card gate run on its statement and proof interval before the card was written.\par\medskip
\noindent\textit{Allowed standard prerequisites (editorial).}\begin{itemize}
\item The characterization of flatness at a point of a morphism of algebraic spaces through local rings. Invoked inside the counted proof (source lines 5536-5545) as a cross-chapter pointer. Supplies the local ring criterion in the scheme-free setting.
\item The description of the local rings of an algebraic stack through a smooth presentation. Used inside the counted proof (source lines 5527-5535). Compares the local rings of the stack with those of the presentation.
\item The statement that flatness is local on the source and target in the smooth topology. Invoked at the end of the counted proof (source lines 5546-5557). Allows the property to be moved between the stack and its presentation.
\end{itemize}\medskip
\ifdefined\hypertarget\hypertarget{sourcepointer-1}{}\fi
\noindent\textit{Source reference.} Original lemma lemma-base-change-flat: sources/stacks-algebraic-stacks-morphisms/pinned-original-source.tex, source lines 5376--5376. Navigation pointer only: the counted text refers to the Stacks item lemma-base-change-flat, which is an existing labeled result of this same chapter and is not part of this card. The pointer carries the target's own label as its number and its definition line as its span; its content is not reproduced here.\par\medskip
\ifdefined\hypertarget\hypertarget{sourcepointer-2}{}\fi
\noindent\textit{Source reference.} Original reference spaces-morphisms-lemma-base-change-module-flat: sources/stacks-algebraic-stacks-morphisms/pinned-original-source.tex, source lines 5428--5428. Navigation pointer only: the counted text refers to the Stacks item spaces-morphisms-lemma-base-change-module-flat, which belongs to another chapter of the project and is not part of this card. The pointer carries the target's own label as its number because this chapter's staged snapshot does not contain the other chapter's tag table; its content is not reproduced here.\par\medskip
\ifdefined\hypertarget\hypertarget{sourcepointer-3}{}\fi
\noindent\textit{Source reference.} Original reference spaces-morphisms-lemma-composition-module-flat: sources/stacks-algebraic-stacks-morphisms/pinned-original-source.tex, source lines 5550--5550. Navigation pointer only: the counted text refers to the Stacks item spaces-morphisms-lemma-composition-module-flat, which belongs to another chapter of the project and is not part of this card. The pointer carries the target's own label as its number because this chapter's staged snapshot does not contain the other chapter's tag table; its content is not reproduced here.\par\medskip
\ifdefined\hypertarget\hypertarget{sourcepointer-4}{}\fi
\noindent\textit{Source reference.} Original reference spaces-morphisms-lemma-flat-permanence: sources/stacks-algebraic-stacks-morphisms/pinned-original-source.tex, source lines 5451--5451. Navigation pointer only: the counted text refers to the Stacks item spaces-morphisms-lemma-flat-permanence, which belongs to another chapter of the project and is not part of this card. The pointer carries the target's own label as its number because this chapter's staged snapshot does not contain the other chapter's tag table; its content is not reproduced here.\par\medskip
\ifdefined\hypertarget\hypertarget{sourcepointer-5}{}\fi
\noindent\textit{Source reference.} Original reference stacks-properties-lemma-points-cartesian: sources/stacks-algebraic-stacks-morphisms/pinned-original-source.tex, source lines 2126--2126. Navigation pointer only: the counted text refers to the Stacks item stacks-properties-lemma-points-cartesian, which belongs to another chapter of the project and is not part of this card. The pointer carries the target's own label as its number because this chapter's staged snapshot does not contain the other chapter's tag table; its content is not reproduced here.\par\medskip
\begin{lemma}
\label{lemma-flat-at-point}
Let $f : \mathcal{X} \to \mathcal{Y}$ be a morphism of algebraic stacks.
Let $x \in |\mathcal{X}|$. Consider commutative diagrams
$$
\vcenter{
\xymatrix{
U \ar[d]_a \ar[r]_h & V \ar[d]^b \\
\mathcal{X} \ar[r]^f & \mathcal{Y}
}
}
\quad\text{with points}
\vcenter{
\xymatrix{
u \in |U| \ar[d] \\
x \in |\mathcal{X}|
}
}
$$
where $U$ and $V$ are algebraic spaces, $b$ is flat, and
$(a, h) : U \to \mathcal{X} \times_\mathcal{Y} V$
is flat. The following are equivalent
\begin{enumerate}
\item $h$ is flat at $u$ for one diagram as above,
\item $h$ is flat at $u$ for every diagram as above.
\end{enumerate}
\end{lemma}

\begin{proof}
Suppose we are given a second diagram $U', V', u', a', b', h'$ as
in the lemma. Then we can consider
$$
\xymatrix{
U \ar[d] & U \times_\mathcal{X} U' \ar[l] \ar[d] \ar[r] & U' \ar[d] \\
V & V \times_\mathcal{Y} V' \ar[l] \ar[r] & V'
}
$$
By Properties of Stacks, Lemma \ref{stacks-properties-lemma-points-cartesian}
there is a point $u'' \in |U \times_\mathcal{X} U'|$ mapping
to $u$ and $u'$. If $h$ is flat at $u$, then the base change
$U \times_V (V \times_\mathcal{Y} V') \to V \times_\mathcal{Y} V'$
is flat at any point over $u$, see
Morphisms of Spaces, Lemma \ref{spaces-morphisms-lemma-base-change-module-flat}.
On the other hand, the morphism
$$
U \times_\mathcal{X} U' \to
U \times_\mathcal{X} (\mathcal{X} \times_\mathcal{Y} V') =
U \times_\mathcal{Y} V' =
U \times_V (V \times_\mathcal{Y} V')
$$
is flat as a base change of $(a', h')$, see Lemma \ref{lemma-base-change-flat}.
Composing and using
Morphisms of Spaces, Lemma \ref{spaces-morphisms-lemma-composition-module-flat}
we conclude that $U \times_\mathcal{X} U' \to V \times_\mathcal{Y} V'$
is flat at $u''$. Then we can use composition by the flat map
$V \times_\mathcal{Y} V' \to V'$ to conclude that
$U \times_\mathcal{X} U' \to V'$ is flat at $u''$.
Finally, since $U \times_\mathcal{X} U' \to U'$ is flat
at $u''$ and $u''$ maps to $u'$ we conclude that
$U' \to V'$ is flat at $u'$ by
Morphisms of Spaces, Lemma \ref{spaces-morphisms-lemma-flat-permanence}.
\end{proof}

\end{document}
